\section{Grafos}
\begin{minted}{cpp}
struct edge{
    int from, to;
    ll w;
    const bool operator<(const edge &b)const{
        return w > b.w;
    }
};
struct pos{
    int from;
    ll c;
    const bool operator<(const pos &b)const{
        return c > b.c;
    }
};
\end{minted}

\vspace{-\baselineskip}

\subsection{Caminos mínimos}
\textbf{Dijkstra.} Complejidad: Tiempo $O(|E| \log |V|)$ - Memoria extra $O(|E|)$.
\begin{minted}{cpp}
ll dijkstra(int a, int b, vector<edge> graph[]){
    ll dist[MAXN];
    bool vis[MAXN] = {};
    fill(dist, dist + MAXN, LLONG_MAX);
    priority_queue<pos> q;
    q.push(pos{a, 0}); dist[a] = 0;
    while(sz(q)){
        pos cur = q.top(); q.pop();
        if(vis[cur.from]) continue;
        vis[cur.from] = true;
        for(edge &e : graph[cur.from]){
            if(dist[e.to] <= dist[e.from] + e.w) continue;
            dist[e.to] = dist[e.from] + e.w;
            q.push(pos{e.to, dist[e.to]});
        }
    } return dist[b];
}
\end{minted}



\textbf{Bellman-Ford.} Complejidad: $O(\abs{V}\abs{E})$.
\begin{minted}{cpp}
vi bellman_ford(int s, int n, vector<edge> &edges, bool cycles = false){
    vi d(n, (cycles ? 0 : INT_MAX));
    d[s] = 0;
    vi P(n, -1); /// Predecesor
    for(int i = 0; i < n - 1; ++i)
    for(edge &e : edges){
        if(d[e.from] == INT_MAX) continue;
        if(d[e.to] > d[e.from] + e.w){
            d[e.to] = d[e.from] + e.w;
            P[e.to] = e.from;
        }
    }
    int last_relax = -1;
    for(edge &e : edges){
        if(d[e.from] == INT_MAX) continue;
        if(d[e.to] > d[e.from] + e.w){
            d[e.to] = d[e.from] + e.w;
            P[e.to] = e.from;
            last_relax = e.to;
        }
    }
    if(last_relax == -1) return d;
    return {}; /// VACIO
}
\end{minted}



\textbf{Floyd-Warshall.} Complejidad: $O\pa{\abs{V}^3}$.
\begin{minted}{cpp}
vvi floyd_warshall(int n){
    const int INF = INT_MAX; // para evitar overflow después
    vvi d(n, vi(n, INF));
    /// aqui inicializa con la lista/matriz de adyacencia
    /// luego calcula la dp
    for(int k = 0; k < n; ++k){
        for(int i = 0; i < n; ++i){
            for(int j = 0; j < n; ++j){
                if(d[i][k]==INF || d[k][j]==INF)
                    continue;
                if(d[i][j] > d[i][k] + d[k][j]) d[i][j] = d[i][k] + d[k][j];
            }
        }
    } return d;
}
\end{minted}



\textbf{Johnson's algorithm.} Complejidad: $O(\abs{V}\abs{E}\log\abs{V})$. Sea $p: V \to \R$ una función potencial del grafo. El algoritmo es como sigue:
\begin{enumerate}[1.]
    \item Hacemos una transformación en el grafo cambiando los pesos $w$ a $w'(u, v) = w(u, v) + p(u) - p(v)$.
    \item Calculamos la distancia mínima $d': V \times V \to \R$ desde cada nodo a todos los demás con Dijkstra.
    \item Finalmente, la distancia mínima de $u$ a $v$ en el grafo original es $d(u, v) = d'(u, v) - p(u) + p(v)$.
\end{enumerate}
La función potencial $p$ puede ser cualquiera. Usando Bellman-Ford se puede calcular el potencial $p(u)$ como el camino más corto que termina (o empieza) en $u$.


\subsection{Árboles}
\textbf{Prim.} Complejidad: Tiempo $O(|E| \log |V|)$. \mintinline{cpp}{eCost[MAXN]} es el arreglo de costos mínimos de cada nodo para incluirlo en el MST.
\begin{minted}{cpp}
ll prim(vector<edge> graph[]){
    ll e_cost[MAXN], ans = 0;
    bool vis[MAXN] = {};
    fill(e_cost, e_cost + MAXN, LLONG_MAX);
    priority_queue<edge> q;
    q.push(edge{1, 1, 0});
    while(sz(q)){
        int node = q.top().to;
        ll w = q.top().w; q.pop();
        if(vis[node]) continue; vis[node] = true;
        for(edge &e : graph[node]){
            if(vis[e.to] || e_cost[e.to] <= e.w) continue;
            e_cost[e.to] = e.w;
            q.push(e);
        } ans += w;
    } return ans;
}
\end{minted}



\textbf{Kruskal.} Complejidad: Tiempo $O\pa{|E|\log |E|}$.
\begin{minted}{cpp}
ll kruskal(vector<edge> &edges, int n){
    sort(all(edges));
    dsu mset(n);
    ll res = 0;
    for(edge &e : edges){
        if(mset.root(e.from) == mset.root(e.to)) continue;
        mset.join(e.from, e.to);
        res += e.w;
    }
    return res;
}
\end{minted}



\textbf{Boruvka.} Complejidad: Tiempo $O\pa{|E|\log |V|}$. $|V| = n$. \mintinline{cpp}{dsu.join()} devuelve \mintinline{cpp}{true} si la unión se llevó a cabo o \mintinline{cpp}{false} en otro caso.
\begin{minted}{cpp}
ll boruvka(vector<edge> &edges, int n){
    dsu mset(n);
    int min_edge[n];
    ll res = 0;
    while(mset.cnt_comp > 1){
        fill(min_edge, min_edge + n, -1);
        for(int i = 0; i < sz(edges); ++i){
            int u = mset.root(edges[i].from);
            int v = mset.root(edges[i].to);
            if(u == v) continue;
            if(min_edge[u] == -1 || edges[i].w < edges[min_edge[u]].w) min_edge[u] = i;
            if(min_edge[v] == -1 || edges[i].w < edges[min_edge[v]].w) min_edge[v] = i;
        }
        for(int i = 0; i < n; ++i){
            int idx_e = min_edge[i];
            if(idx_e == -1) continue;
            res += mset.join(edges[idx_e].from, edges[idx_e].to) * edges[idx_e].w;
        }
    } return res;
}
\end{minted}



\textbf{MST dirigido.} Complejidad: $O(|E|\log|V|)$. AGREGAR PEQUEÑA DESCRIPCIÓN.
\begin{minted}{cpp}
/// MEJORAR ESTA COSA, SOLO LO COPIE Y PEGUÉ porcuestionesdetiempo
const int N = 3e5 + 9;
const ll inf = 1e18;
template<typename T> struct PQ {
    ll sum = 0;
    priority_queue<T, vector<T>, greater<T>> Q;
    void push(T x) { x.w -= sum; Q.push(x); }
    T pop() { auto ans = Q.top(); Q.pop(); ans.w += sum; return ans; }
    int size() { return sz(Q); }
    void add(ll x) { sum += x; }
    void merge(PQ &x){
        if (size() < sz(x)){
            swap(sum, x.sum);
            Q.swap(x.Q);
        }
        while(sz(x)){
            auto tmp = x.pop();
            tmp.w -= sum;
            Q.push(tmp);
        }
    }
};
struct edge {
    int u, v; ll w;
    bool operator > (const edge &rhs) const { return w > rhs.w; }
};
struct DSU {
    vi par;
    DSU (int n) : par(n, -1) {}
    int root(int i) { return par[i] < 0 ? i : par[i] = root(par[i]); }
    void set_par(int c, int p) { par[c] = p; }
};
// returns parents of each vertex
// each edge should be distinct
// it assumes that a solution exists (all vertices are reachable from root)
// 0 indexed
// Takes ~300ms for n = 2e5
vi DMST(int n, int root, const vector<edge> &edges) {
    vi u(2 * n - 1, -1), par(2 * n - 1, -1);
    edge par_edge[2 * n - 1];
    vi child[2 * n - 1];
    PQ<edge> Q[2 * n - 1];
    for(auto e : edges) Q[e.v].push(e);
    for(int i = 0; i < n; i++)
        Q[(i+1) % n].push({i, (i+1) % n, inf});
    int super = n;
    DSU dsu(2 * n - 1);
    int head = 0;
    while(sz(Q[head])) {
        auto x = Q[head].pop();
        int nxt_root = dsu.root(x.u);
        if (nxt_root == head) continue;
        u[head] = nxt_root;
        par_edge[head] = x;
        if (u[nxt_root] == -1) head = nxt_root;
        else {
            int j = nxt_root;
            do {
                Q[j].add(-par_edge[j].w);
                //Q[super].merge(move(Q[j]));
                Q[super].merge(Q[j]);
                assert(u[j] != -1);
                child[super].pb(j);
                j = dsu.root(u[j]);
            } while (j != nxt_root);
            for(auto u : child[super]) par[u] = super, dsu.set_par(u, super);
            head = super++;
        }
    }
    vi res(2 * n - 1, -1);
    queue<int> q; q.push(root);
    while (sz(q)) {
        int u = q.front();
        q.pop();
        while (par[u] != -1) {
            for (auto v : child[par[u]]) {
                if (v != u) {
                    res[par_edge[v].v] = par_edge[v].u;
                    q.push(par_edge[v].v);
                    par[v] = -1;
                }
            }
            u = par[u];
        }
    }
    res[root] = root; res.resize(n);
    return res;
}
int main() {
    ios_base::sync_with_stdio(0); cin.tie(0);
    int n, m, root; cin >> n >> m >> root;
    vector<edge> edges(m);
    for(auto &e : edges) cin >> e.u >> e.v >> e.w;
    auto res = DMST(n, root, edges);
    unordered_map<int, int> W[n];
    for (auto u : edges) W[u.v][u.u] = u.w;
    ll ans = 0;
    for(int i = 0; i < n; i++) if(i != root)
        ans += W[i][res[i]];
    cout << ans << '\n';
    for(auto x : res) cout << x << ' ';
    cout << '\n';
}
\end{minted}



\textbf{Block-Cut Tree} Complejidad: $O(n)$?
\begin{minted}{cpp}
/// MEJORAR ESTA COSA, SOLO LO COPIE Y PEGUÉ porcuestionesdetiempo
vvi comps;//nodes in each component
vvi bcc;//nodes to components that it belong
vi st;//internal stack
vi low, num;
int id;
int sz;
void dfs_biconnected(vvi &g, int u, int pre){
  low[u] = num[u] = ++id;
  st.pb(u);
  for(auto v: g[u]) {
    if(!num[v]) {
      dfs_biconnected(g, v, u);
      low[u] = min(low[u], low[v]);
      if(low[v] >= num[u]) {
        int x;
        comps.pb(vi());
        do {
          x = st.back();
          st.pop_back();
          bcc[x].pb(sz);
          comps.back().pb(x);
        } while(x ^ v);
        bcc[u].pb(sz);
        comps.back().pb(u);
        sz++;
      }
    } else if(v != pre) low[u] = min(low[u], num[v]);
  }
}
vi utoubt;//its component or its AP index
vb uisart;//u is AP?
vvi bt;//block cut tree
void generateBlockCutTree(vvi &g){
    int n = sz(g);
    sz = id = 0;
    bcc.resize(0); low.resize(0); num.resize(0);
    bcc.resize(n); low.resize(n); num.resize(n);
    dfs_biconnected(g, 0, 0);
    bt.resize(sz);
    utoubt.resize(0); utoubt.resize(n);
    uisart.resize(0); uisart.resize(n);
    for(int u = 0; u < n; u++){
        if(sz(bcc[u]) > 1) { //Articulation
            utoubt[u] = sz++;
            uisart[u] = true;
            bt.pb(vi());
            for(auto v: bcc[u]){
                bt[utoubt[u]].pb(v);
                bt[v].pb(utoubt[u]);
            }
        } else //Not articulation point
            utoubt[u] = bcc[u][0];
    }
}
\end{minted}



\textbf{LCA.} Complejidad: Tiempo de preproceso $O(|V|\log |V|)$. Tiempo de LCA y $n$-ésimo ancestro $O(\log |V|)$. $1$-index.
\begin{minted}{cpp}
void precalc(int node, int p = 0, int d = 1){
    depth[node] = d;
    P[0][node] = p;
    for(int k = 1; k <= LOGN; ++k)
        P[k][node] = P[k - 1][P[k - 1][node]];
    for(int child : tree[node]) if(p != child)
        precalc(child, node, d + 1);
}
int LCA(int a, int b){
    if(depth[b] < depth[a]) swap(a, b);
    int dif = depth[b] - depth[a];
    for(int k = LOGN; 0 <= k; --k)
        if(is_on(dif, k)) b = P[k][b];
    if(a == b) return a;
    for(int k = LOGN; 0 <= k; --k){
        if(P[k][a] != P[k][b]){
            a = P[k][a];
            b = P[k][b];
        }
    }
    return P[0][a];
}
int nth_ancestor(int u, int n){
    for(int k = LOGN; 0 <= k; --k)
        if(is_on(n, k)) u = P[k][u];
    return u;
}
\end{minted}



\textbf{Sack.} Complejidad: Tiempo $O(\abs{V}\log \abs{V})$. $1$-index.
\begin{minted}{cpp}
void precalc(int node, int p = 0){
    subtree_size[node] = 1;
    depth[node] = depth[p] + 1;
    for(int v : tree[node]){
        if(v == p) continue;
        precalc(v, node);
        subtree_size[node] += subtree_size[v];
    }
}
void add(int node, int x, int p = 0){
    /// add node here
    /// add subtree
    for(int v: tree[node])
        if(v != p && !big[v])
            add(v, x, node);
}
void dfs(int node, bool keep, int p = 0){
    int maxi = -1, big_child = -1;
    for(int v : tree[node])///Searchfor big_child
       if(v != p && subtree_size[v] > maxi)
          maxi = subtree_size[v], big_child = v;
    for(int v : tree[node])
        if(v != p && v != big_child)
            dfs(v, false, node);  /// run a dfs on small childs and clear them
    if(big_child != -1)
        dfs(big_child, true, node), big[big_child] = 1;  /// big_child marked as big and not cleared
    add(node, 1, p);
    /// answer queries here
    if(big_child != -1) big[big_child] = 0;
    if(!keep) add(node, -1, p);
}
\end{minted}



\subsection{Máximo flujo}

\textbf{Algunos problemas de flujos}
\begin{itemize}
    \item \textit{Maximum Weight Closure}. Sea $N_1$ una clausura de $G$ y $N_2 = V \setminus N_1$, tenemos que $w\pa{N_1} = \sum_{i \in N_1^+} w_i - \sum_{i \in N_1^-} |w_i|$ y $Cap.Corte = \sum_{i \in N_2^+} w_i + \sum_{i \in N_1^-} |w_i|$. Entonces
    $$Cap.Corte + w(N_1) = \sum_{i \in N_1^+} w_i + \sum_{i \in N_2^+} w_i.$$
    \item \textit{Mínima cobertura de vértices}. En grafos generales es NP-Completo. En grafos bipartitos el máximo emparejamiento es igual al numero de vertices en la mínima cobertura. Para el problema con pesos en los nodos, unimos $s$ a todos los nodos en $L$ con capacidad igual al peso de cada nodo, unimos los nodos de $R$ a $t$ de la misma manera y unimos los nodos de $L$ a $R$ con capacidad infinita. El máximo flujo es el peso mínimo de la mínima cobertura.
    \item \textit{Máximo conjunto independiente}. Cualquier conjunto independiente es el complemento de alguna cobertura de vértices.
    \item \textit{Mínimo cubrimiento de caminos independientes}. En grafos generales es NP-hard. En DAG's duplicamos los nodos en un lado IN y un lado OUT. Conectamos $s$ al lado OUT y el lado IN a $t$. Las aristas del DAG las agregamos del lado OUT al lado IN. Sea $M$ el máximo emparejamiento de la red anterior, entonces el mínimo cubrimiento es $\abs{V} - M$.
    \item \textit{Mínimo cubrimiento de caminos NO necesariamente independientes}. En grafos generales es NP-hard. En DAG's transformamos el DAG a su clausura transitiva y aplicamos el problema anterior.
    \item \textit{Teorema de Menger}. Sea $G$ un grafo finito (sin pesos) con vértices $s$, $t$ distintos. El corte mínimo (por aristas) entre $s$ y $t$ es igual a la cantidad máxima de caminos independientes por aristas de $s$ a $t$. Si $s$ y $t$ no son adyacentes, el corte mínimo por vértices entre $s$ y $t$ es igual a la cantidad máxima de caminos independientes por vértices de $s$ a $t$.
    \item \textit{Teorema de Mirsky}. En todo POSET, el tamaño de la cadena de mayor tamaño es igual al número de anticadenas necesarias para cubrir todos los elementos del conjunto.
    \item \textit{Teorema de Dilworth}. En todo POSET, el tamaño de la anticadena de mayor tamaño es igual al número de cadenas necesarias para cubrir todos los elementos del conjunto.
    \item \textit{Teorema de Hall}. Un grafo bipartito con subconjuntos $L$ y $R$ tiene un emparejamiento de tamaño $|L|$ si y sólo si para todo subconjunto $W$ de $L$, se cumple que $|W| \leq |N_G(W)|$, donde $N_G(W)$ es el conjunto de vértices vecinos de alguno de los vértices en $W$.
\end{itemize}


\textbf{Edmonds-Karp}
Complejidad: Ford-Fulkerson $O(|E| \cdot F)$, Edmonds-Karp $O(\abs{V}\abs{E}^2)$.
\begin{minted}{cpp}
struct edge {
    int from, to;
    ll w, c, f;
    // weight, capacity, flow
};
class ford_fulkerson {
public:
    ford_fulkerson (vector<vector<edge>> &graph) : graph(graph){}
    ll get_max_flow(int s, int t){
        init();
        ll f = 0;
        while(find_and_update(s, t, f)){}
        return f;
    }
    vi get_st_cut(const int &s){
        bool vis[sz(graph)] = {};
        vi S;
        queue<int> q;
        q.push(s);
        S.pb(s);
        vis[s] = true;
        while(sz(q)){
            int u = q.front(); q.pop();
            for(int eI : edge_indexes[u]){
                if(edges[eI].c > edges[eI].f && !vis[edges[eI].to]){
                    q.push(edges[eI].to);
                    vis[edges[eI].to] = true;
                    S.pb(edges[eI].to);
                }
            }
        }
        return S;
    }
    vector<vector<edge>> get_residual_graph(){
        vector<vector<edge>> residual(sz(graph));
        for(int i=0; i<sz(edges); i+=2){
            const edge& e = edges[i];
            if(e.c > 0){
                residual[e.from].pb({e.from, e.to, e.w, e.c - e.f, e.f});
                residual[e.to].pb({e.to, e.from, -e.w, e.f, -e.f});
            }
        } return residual;
    }
private:
    vector<vector<edge>> graph;
    vector<edge> edges;
    vvi edge_indexes;
    void init(){
      edges.clear();
      edge_indexes.clear();
      edge_indexes.resize(sz(graph));
      for(int u = 0; u < sz(graph); u++){
         for(edge &e : graph[u]){
           edges.pb({u, e.to, e.w, e.c, 0});
           edges.pb({e.to, u, -e.w, 0, 0});
           edge_indexes[u].pb(sz(edges)-2);
           edge_indexes[e.to].pb(sz(edges)-1);
         }
      }
    }
    bool find_and_update(int s, int t, ll &flow){
        queue<int> q;
        // Desde donde llego y con que arista
        vpii from(sz(graph), {-1, -1});
        q.push(s);
        from[s] = {s, -1};
        bool found = 0;
        while(sz(q) && !found){
            int u = q.front(); q.pop();
            for(int eI : edge_indexes[u]){
                if(edges[eI].c > edges[eI].f &&
                  from[edges[eI].to].fi== -1){
                    from[edges[eI].to] = {u, eI};
                    q.push(edges[eI].to);
                    if(edges[eI].to == t)found=1;
                }
            }
        }
        if(!found) return false;
        ll u_flow = LLONG_MAX;
        int cur = t;
        while(cur != s) {
            u_flow = min(u_flow, edges[from[cur].se].c - edges[from[cur].se].f);
            cur = from[cur].fi;
        }
        cur = t;
        while(cur != s){
            edges[from[cur].se].f += u_flow;
            edges[from[cur].se^1].f -= u_flow;
            cur = from[cur].fi;
        }
        flow += u_flow;
        return 1;
    }
};
\end{minted}



\textbf{Dinic.} Complejidad: $O(\abs{V}^2\abs{E})$.
\begin{minted}{cpp}
const int MAXV = 32767; /// 2^15 - 1
template<class T = ll> struct dinic{
    dinic(short V){this->V = V;}
    const static bool SCALING = 1;
    bool sorted = 0;
    short s, t, V;
    int lim = 1; /// Para escalado
    const T INF = numeric_limits<T>::max();
    short level[MAXV]; /// distancia desde s
    short ptr[MAXV]; /// arista que va explorando
    struct edge{
        short to, rev;
        T cap, flow, mcap;
        bool operator<(const edge &b)const{return mcap > b.mcap;}
    };
    vector<edge> adj[MAXV];
    vi adj_cur[MAXV]; /// aristas del grafo de nivel
    void add_edge(short u, short v, T cap, bool is_directed = 1){
        if(u == v) return;
        T add = (is_directed ? 0 : cap);
        adj[u].pb({v, sz(adj[v]), cap, 0 ,cap+add});
        adj[v].pb({u, sz(adj[u])-1, add, 0, cap+add});
    }
    void mysort(){
        if(sorted) return;
        sorted = 1;
        for(int i = 0; i < V; ++i){
            sort(all(adj[i]));
            for(int j=0; j<sz(adj[i]); ++j)
            adj[adj[i][j].to][adj[i][j].rev].rev = j;
        }
    }
    bool bfs(){
        for(int i = 0; i < V; ++i){
            adj_cur[i].clear();
            adj_cur[i].reserve(sz(adj[i]));
        }
        queue<short> q;
        q.push(s);
        fill(level, level + V, -1);
        level[s] = 0;
        while(sz(q)){
            short u = q.front(); q.pop();
            if(u == t) return 1;
            for(int i=0; i < sz(adj[u]); ++i){
                edge &e = adj[u][i];
                if(e.mcap < lim) break;
                if(level[e.to] == -1 && e.cap - e.flow >= lim){
                    level[e.to] = level[u] + 1;
                    adj_cur[u].pb(i);
                    q.push(e.to);
                } else if(level[e.to]==level[u]+1 && e.cap - e.flow >= lim){
                    adj_cur[u].pb(i);
                }
            }
        } return 0;
    }
    T dfs(short u, T flow, vector<short> &S, bool save = 0){
        if(save) S.pb(u);
        if(u == t) return flow;
        for(; ptr[u] < sz(adj_cur[u]); ++ptr[u]){
            edge &e = adj[u][adj_cur[u][ptr[u]]];
            if(T pushed = dfs(e.to, min(flow, e.cap - e.flow), S, save)){
                e.flow += pushed;
                adj[e.to][e.rev].flow -= pushed;
                if(e.cap-e.flow < lim) ptr[u]++;
                return pushed;
            }
        } return 0;
    }
    ll get_max_flow(short source, short sink){
        s = source; t = sink;
        mysort();
        vector<short> S;
        ll flow = 0;
        lim = SCALING ? (1<<30) : 1;
        for(; 0 < lim; lim >>= 1){
            while(bfs()){
                memset(ptr, 0, sizeof(ptr));
                while(T pushed = dfs(s, INF, S))
                    flow += pushed;
            }
        } return flow;
    }
    vector<short> get_st_cut(){
        vector<short> S;
        memset(ptr, 0, sizeof(ptr));
        dfs(s, INF, S, true);
        return S;
    }
};
\end{minted}



\textbf{MCMF.} Complejidad: $O(\min\{|E|^2|V|\log |V|, \; F|E|\log |V|\}).$
\begin{minted}{cpp}
const int MAXV = (1 << 15) - 1;
template<class T = ll> struct mcmf{
    mcmf(short V){this->V = V;}
    short s, t, V;
    const T INF = numeric_limits<T>::max();
    vector<T> p;
    struct edge{
        short to, rev;
        T cap, flow, mcap, w;
        bool operator<(const edge &b)const{return mcap > b.mcap;}
    };
    vector<edge> adj[MAXV];
    void add_edge(short u, short v, T cap, T w, bool is_directed = 1){
        if(u == v) return;
        T add = (is_directed ? 0 : cap);
        adj[u].pb({v, sz(adj[v]), cap, 0, cap + add, w});
        adj[v].pb({u, sz(adj[u]) - 1, add, 0, cap + add, -w});
    }
    struct pos{
        short from;
        T c;
        const bool operator<(const pos &b)const{return c > b.c;}
    };
    vector<T> bellman_ford(){
        vector<T> d(V);
        for(int i = 0; i < V - 1; ++i){
            for(int u = 0; u < V; ++u){
                for(edge &e : adj[u]){
                    if(d[e.to] > d[u] + e.w && e.cap > 0){
                        d[e.to] = d[u] + e.w;
                    }
                }
            }
        } return d;
    }
    pair<T, T> dijkstra(const T MAX_FLOW){
        vector<T> d(V, INF);
        vi P(V, -1), P_e(V, -1);
        priority_queue<pos> q;
        q.push({s, 0});
        d[s] = 0;
        while(sz(q)){
            pos act = q.top(); q.pop();
            if(act.c != d[act.from]) continue;
            for(int j=0; j<sz(adj[act.from]); ++j){
                edge &e = adj[act.from][j];
                if(e.cap - e.flow <= 0) continue;
                T _w = e.w+p[act.from] - p[e.to];
                if(d[e.to] <= d[act.from] + _w)
                    continue;
                d[e.to] = d[act.from] + _w;
                q.push(pos{e.to, d[e.to]});
                P[e.to] = act.from;
                P_e[e.to] = j;
            }
        }
        for(int i = 0; i < V; ++i) if(d[i] < INF)
            d[i] += -p[s] + p[i];
        for(int i = 0; i < V; ++i) if(d[i] < INF)
            p[i] = d[i];
        if(P[t] == -1) return {0, 0};
        T flow = MAX_FLOW;
        int cur_node = t;
        while(cur_node != s){
            flow = min(flow, adj[ P[cur_node] ][ P_e[cur_node] ].cap - adj[ P[cur_node] ][ P_e[cur_node] ].flow);
            cur_node = P[cur_node];
        }
        T new_cost = 0;
        cur_node = t;
        while(cur_node != s){
            adj[ P[cur_node] ][ P_e[cur_node] ].flow += flow;
            new_cost += adj[ P[cur_node] ][ P_e[cur_node] ].w * flow;
            adj[cur_node][adj[ P[cur_node] ][ P_e[cur_node] ].rev ].flow -= flow;
            cur_node = P[cur_node];
        }
        return {flow, new_cost};
    }
    pair<T, T> get_max_flow(short source, short sink, const T MAX_FLOW = 1e8){
        s = source;
        t = sink;
        p = bellman_ford();
        T flow = 0, cost = 0;
        while(flow < MAX_FLOW){
            pair<T, T> pushed = dijkstra(MAX_FLOW - flow);
            if(!pushed.fi) break;
            flow += pushed.fi;
            cost += pushed.se;
        }
        return {flow, cost};
    }
};
\end{minted}



\subsection{SCC}
\textbf{Kosajaru}
Complejidad: Tiempo $O(n)$.
\begin{minted}{cpp}
void dfs(int u, vi graph[], bool vis[], vi &topo_ord){
    if(vis[u]) return;
    vis[u] = true;
    for(int v : graph[u])dfs(v, graph, vis, topo_ord);
    topo_ord.pb(u);
}
void assign_scc(int u, vi inv_graph[], bool vis[], vi &scc, const int id){
    if(vis[u]) return;
    vis[u] = true;
    scc[u] = id;
    for(int v : inv_graph[u])
        assign_scc(v, inv_graph, vis, scc, id);
}
pair<int, vi> kosajaru(int n, vi graph[], vi inv_graph[]){
    bool vis[n] = {};
    vi topo_ord;
    for(int i=0; i<n; ++i) dfs(i,graph,vis,topo_ord);
    reverse(all(topo_ord));
    memset(vis, 0, sizeof(vis));
    vi scc(n);
    int id = 0;
    for(int u : topo_ord) if(!vis[u])
        assign_scc(u, inv_graph, vis, scc, id++);
    return {id, scc};
}
pair<vvi, vvi> build_scc_graph(int n, vi graph[], int max_scc, const vi &scc){
    vvi scc_graph(max_scc), inv_scc_graph(max_scc);
    for(int u = 0; u < n; ++u)
    for(int v : graph[u])
    if(scc[u] != scc[v])
        scc_graph[scc[u]].pb(scc[v]);
    for(int u = 0; u < n_scc; ++u){
        sort(all(scc_graph[u]));
        auto it = unique(all(scc_graph[u]));
        scc_graph[u].resize(it - scc_graph[u].begin());
        for(int v : scc_graph[u])
            inv_scc_graph[v].pb(u);
    } return {scc_graph, inv_scc_graph};
}
\end{minted}



\subsection{2-Sat}
Complejidad: Tiempo en responder $O(n)$.
\begin{minted}{cpp}
struct two_sat{
    int n;
    vvi graph, inv_graph;
    vi scc, ans;
    vector<bool> vis;
    two_sat(){}
    two_sat(int _n){
        n = _n;
        graph.resize(2 * n);
        inv_graph.resize(2 * n);
        scc.resize(2 * n);
        vis.resize(2 * n);
        ans.resize(n);
    }
    void add_edge(int u, int v){
        graph[u].pb(v);
        inv_graph[v].pb(u);
    }
    /// al menos una es verdadera
    void add_or(int p, bool val_p, int q, bool val_q){
        add_edge(p+(val_p? n:0), q+(val_q? 0:n));
        add_edge(q+(val_q? n:0), p+(val_p? 0:n));
    }
    /// exactamente una es verdadera
    void add_xor(int p, bool val_p, int q, bool val_q){
        add_or(p, val_p, q, val_q);
        add_or(p, !val_p, q, !val_q);
    }
    /// p y q tienen el mismo valor
    void add_and(int p, bool val_p, int q, bool val_q){
        add_xor(p, !val_p, q, val_q);
    }
    /// Kosajaru
    void dfs(int node, vi &topo_ord){...}
    void assign_scc(int node, const int id){...}
    /// construye respuesta
    bool build_ans(){
        fill(all(vis), false);
        vi topo_ord;
        for(int i=0; i<2*n; ++i)dfs(i, topo_ord);
        fill(all(vis), false);
        reverse(all(topo_ord));
        int id = 0;
        for(int u : topo_ord) if(!vis[u])
            assign_scc(u, id++);
        for(int i = 0; i < n; ++i){
            if(scc[i] == scc[i + n]) return 0;
            ans[i] = (scc[i]<scc[i + n] ? 0:1);
        } return 1;
    }
};
\end{minted}